\documentclass[12pt]{scrartcl}
\usepackage[a4paper,margin=3.0cm]{geometry}
\usepackage{../common/rlstudy}
\linespread{1.15}\selectfont

\title{Extended reading --- Week 02: Bellman Equations and Dynamic Programming}
\date{\today}
\begin{document}
\maketitle

\section*{How to use this list}
Read these five in order. Bellman (1957) and Howard (1960) are the
founding documents; Puterman (1994) is the canonical modern textbook
chapter you should re-read every couple of years; Tsitsiklis \& Van Roy
(1997) explains \emph{why} naive DP-with-function-approximation diverges
and motivates the contraction-preserving variants you will meet in
Week 5 onward; Munos \& Szepesv\'ari (2008) bridges DP to learning by
delivering the first proper finite-sample bounds for approximate
value iteration.

\section{Bellman (1957) --- Dynamic Programming, Chapters 1--3}
\textit{Princeton University Press. Republished by Dover, 2003.}

\paragraph{Summary.} Bellman defines the \emph{principle of optimality}
(``an optimal policy has the property that whatever the initial state and
initial decision are, the remaining decisions must constitute an optimal
policy with regard to the state resulting from the first decision'') and
derives the recursion that bears his name. The book is short, conversational,
and oddly readable for something written in the mid-1950s.

\paragraph{Why it matters for this week.} The principle of optimality is the
\emph{semantic} content of the Bellman optimality equation. Knowing the
equation is one thing; knowing it is a statement about the recursive structure
of optimal substructures is what lets you recognise DP when it shows up in
disguise (the LP formulation, the linear-system formulation in Eq.~3.8 of the
notes, the policy-gradient theorem).

\begin{quote}\itshape
``In place of determining the optimal sequence of decisions from the fixed
state of the system, we wish to determine the optimal decision to be made at
any state of the system. Only if we know the latter do we understand the
intrinsic structure of the solution.''
\end{quote}
The shift from ``sequence of decisions'' to ``decision per state'' is the move
from open-loop control to closed-loop policy --- the conceptual leap that
makes RL possible.

\section{Howard (1960) --- Dynamic Programming and Markov Processes}
\textit{MIT Press / Wiley. Chapters 4--6.}

\paragraph{Summary.} Howard introduces policy iteration (the algorithm that
bears his name) and proves it converges in finitely many steps for finite
MDPs. The proof uses the policy-improvement theorem of \S 3.4 of your notes;
the finiteness comes from the strict improvement at every step combined with
the finite number of deterministic stationary policies.

\paragraph{Why it matters for this week.} Howard's PI is the algorithm with
\emph{strongest convergence properties} in DP: a guaranteed finite number of
sweeps to the exact optimum, no $\epsilon$ tolerance. Read it for the cleanest
example of a structural-induction proof in optimisation.

\begin{quote}\itshape
``A policy that is unimproved by policy improvement is optimal\ldots\ Since
there are only finitely many policies, the algorithm must terminate in finitely
many iterations.''
\end{quote}
The proof is essentially three sentences. Howard packs the argument into a
chapter of motivation, but the core idea is this: each iteration strictly
improves, and there are only $|\Actions|^{|\States|}$ policies, so termination
is forced.

\section{Puterman (1994) --- Markov Decision Processes, Chapters 6--7}
\textit{Wiley.}

\paragraph{Summary.} The definitive modern treatment of finite-MDP DP.
Chapter 6 contains the proof of contraction (your Theorem 3.2), the
existence and uniqueness of $V^\star$, and a polynomial-time bound for
policy iteration. Chapter 7 covers modified policy iteration with the proof
that any $k\ge 1$ inner sweeps suffice for convergence, plus action
elimination (a precursor to Q-learning's targeted updates).

\paragraph{Why it matters for this week.} The textbook proofs in Puterman are
the cleanest you will find, and they recur as templates in every subsequent
DP and ADP analysis. Memorise the inequality
$|\max_a f(a) - \max_a g(a)| \le \max_a |f(a) - g(a)|$ from \S 6.2; you
will use it constantly.

\begin{quote}\itshape
``Modified policy iteration\ldots\ unifies value iteration ($k=1$) and policy
iteration ($k=\infty$) as endpoints of a single family of algorithms. The
convergence rate is geometric for every $k$, with constant depending on $k$
and $\Disc$.''
\end{quote}
Treating VI and PI as endpoints of a single family is the right mental model.
When we move to learning, ``how many off-policy updates per policy improvement
step?'' is the same question.

\section{Tsitsiklis \& Van Roy (1997) --- An analysis of temporal-difference
learning with function approximation}
\textit{IEEE TAC 42(5), 674--690. Available also as MIT LIDS-P-2390.}

\paragraph{Summary.} Proves that TD(0) with \emph{linear} function approximation
and on-policy sampling converges to the projection of $V^\Policy$ onto the
linear subspace, in a weighted sup-norm. Off-policy, with linear FA, the
same algorithm can \emph{diverge} to infinity --- the famous Baird counterexample
is the canonical demonstration. The paper is the reason ``deadly triad'' is
a phrase in RL.

\paragraph{Why it matters for this week.} It is the bridge from exact DP
(this week) to approximate DP (Weeks 5, 6, 7, 10). Knowing that
\emph{contraction is preserved on-policy but not off-policy under linear FA}
explains why DQN needed target networks and Polyak averaging; why TD3 needed
twin critics; why Soft Actor-Critic uses entropy regularisation to keep its
critic in a contraction basin.

\begin{quote}\itshape
``The combination of function approximation, off-policy training, and
bootstrapping --- the \emph{deadly triad} --- is the source of TD divergence.
On-policy linear TD(0) converges; remove any one of the three and divergence
is precluded; have all three and divergence is generically possible.''
\end{quote}
The deadly-triad lens is invaluable: when a deep-RL training run diverges,
the first diagnostic is to identify which of the three is at play and whether
the algorithm is supposed to handle it.

\section{Munos \& Szepesv\'ari (2008) --- Finite-time bounds for fitted value
iteration}
\textit{JMLR 9, 815--857.}

\paragraph{Summary.} The first finite-sample analysis of \emph{approximate}
value iteration with general function approximation. Given a regression
oracle and an exploration distribution $\mu$, the value-iteration error after
$K$ iterations is bounded by a sum of (i) the projection error of $T_\star V_K$
onto the function class, (ii) the concentrability coefficient
$C_\mu$ that measures how well $\mu$ covers the on-policy distribution, and
(iii) a $\gamma^K$ term that decays geometrically. The paper invented the
modern ADP error decomposition you will see in offline-RL papers.

\paragraph{Why it matters for this week.} It tells you what bills come due
when you replace the table of $V$ by a neural net: regression error plus a
distribution-shift coefficient. The $C_\mu$ coefficient is the formal version
of the intuition ``you cannot evaluate $\Policy$ in regions the data never
visits.'' All of offline RL (Weeks 14, 21) is engineering against $C_\mu$.

\begin{quote}\itshape
``The performance loss is bounded by a sum involving the approximation error,
the concentrability of the future-state distribution with respect to the
sampling distribution, and a $\gamma^K$ term that vanishes geometrically.''
\end{quote}
``Concentrability of the future-state distribution'' is the abstract version
of ``do not extrapolate''. It is also the reason behaviour-cloning baselines
beat naive offline-RL on out-of-distribution states.

\section*{Reading exercises}
\begin{enumerate}\setlength\itemsep{4pt}
  \item \textbf{Bellman, the LP formulation.} Read \S 1.3 of Bellman (1957).
        The value function $V^\star$ is the minimum-norm vector $V$ satisfying
        $V \ge T_a V$ for all $a$. Write down the linear program over $V$ and
        identify the dual variables. (Answer: the dual variables are the
        \emph{discounted state-occupancy measure} of the optimal policy.)
  \item \textbf{Howard, finite-time PI bound.} Howard's policy iteration is
        guaranteed to terminate in at most $|\Actions|^{|\States|}$ outer
        iterations. Ye (2011, \emph{Math.\ Oper.\ Res.}) proved a stronger
        bound: PI terminates in $O(|\States|\,|\Actions|/(1-\gamma) \log(\,
        |\States|/(1-\gamma)))$ outer iterations. Sketch the argument; why
        is the bound \emph{strongly polynomial} when $\gamma$ is fixed?
  \item \textbf{Puterman, action elimination.} Read \S 6.7. Implement
        action elimination on a 10x10 GridWorld and compare the per-sweep
        cost vs vanilla VI. By how much does it accelerate convergence at
        $\gamma=0.99$?
  \item \textbf{Tsitsiklis \& Van Roy, Baird's counterexample.} Reproduce
        Baird's 7-state ``star'' MDP with linear features. Run off-policy
        TD(0) with linear FA and a fixed behaviour. Watch it diverge to
        infinity. Now replace TD(0) with \emph{gradient-TD} (Sutton et al.\
        2009) and watch it converge. What does the gradient-TD fix change
        in the update rule?
  \item \textbf{Munos--Szepesv\'ari, concentrability in offline RL.} Read
        Theorem 4 of the paper and apply it to a synthetic offline-RL
        problem where the behaviour distribution is uniform but the optimal
        on-policy distribution is concentrated. What is $C_\mu$? How does
        it scale as the gap between the two distributions grows?
\end{enumerate}

\section*{Connections to later weeks}
\begin{itemize}\setlength\itemsep{2pt}
  \item \emph{Week 4 (TD, SARSA, Q-learning)} replaces the expectation in
        $T_\star$ by a single-sample stochastic gradient; the contraction
        becomes a stochastic-approximation Lyapunov argument.
  \item \emph{Week 5 (DQN)} adds neural-net function approximation and is
        the first place where ``contraction is preserved''
        breaks --- which is why DQN needs target networks and replay.
  \item \emph{Week 7 (policy gradients)} re-derives improvement via the
        \emph{policy gradient theorem}, an explicit closed-form
        version of the improvement step from Howard's PI.
  \item \emph{Week 14 (offline RL)} pays the concentrability bill from
        Munos--Szepesv\'ari in the form of pessimism penalties (CQL, IQL).
\end{itemize}
DP is the spine of every algorithm in the rest of the course. Returning
to this week's contraction proof when a deep-RL method diverges is one of
the most productive things you can do as a practitioner.

\section*{Sourced from}
\verb|research/deeprl_core_sota.md| \S 1 and Sutton \& Barto Chapter 4. For
modern DP-with-LP formulation see Bertsekas, \emph{Dynamic Programming and
Optimal Control} Vol.\ II (Athena Scientific, 4th ed., 2017).

\end{document}
